\section{The observation class and the main results}\label{sec:setup}
The problem is to restore a fixed coercivity target using few permitted
observations. Two distinctions matter: the dimension of a displayed output
space versus its actual rank on the stated domain, and a finite penalty versus
a penalty bounded independently of the lattice size. We keep both explicit.
The proofs concern one positive family of finite operators, rather than a
uniform theorem for arbitrary matrix profiles.

\subsection{Actual finite gradients and the coexact domain}
Fix integers $D\geq3$, $n\geq2$ and $g\geq1$, and write
\[
 X=(\Z/n\Z)^D,\qquad V=n^D,\qquad
 \mathcal H_g=(\R^D\otimes\R^g)^X.
\]
All unqualified inner products and norms use counting measure. Put
$\partial_i f(x)=f(x+e_i)-f(x)$, retaining every positive forward label,
including period-two multiplicities. For $z\in\mathcal H_g$ let
\[
 \En(z)=\sum_{x\in X}\sum_{i,k=1}^D\sum_{\alpha=1}^g
       |\partial_i z_{k,\alpha}(x)|^2.
\]
For a real one-form $v$ define
$(d_1v)_{ij}=\partial_i v_j-\partial_j v_i$ for $i<j$.
The adjoint uses the counting inner products, with each unordered pair $i<j$
counted once. Our domain is
\begin{equation}\label{eq:domain}
 \UU_0=\im d_1^*,\qquad \UU=\UU_0^g\subset\mathcal H_g,
 \qquad d_0:=\dim_\R\UU_0=(D-1)(V-1).
\end{equation}
The dimension identity is proved in Lemma~\ref{lem:fourier}. No artificial
componentwise replacement of this domain is used.

For a positive symmetric $g\times g$ field $B(x)$ set
\begin{equation}\label{eq:energy}
 Q_B(z)=\sum_{x\in X}\sum_{i,k=1}^D
       (\partial_i z_k(x))^T B(x)\partial_i z_k(x).
\end{equation}
The coefficient is evaluated at the actual forward \emph{source} $x$.
The full gradient-stack profile is $\mathcal A(x)=I_{D^2}\otimes B(x)$;
there are $D^2g$ local eigenvalues, while the field has $h=Dg$ components.
This notation does not assert an identity with an interacting or continuum
operator.

\subsection{Canonical observations and normalization}
For each original internal channel $\alpha$, choose a partition
$\mathcal P_\alpha$ of $\{0,\ldots,n-1\}^D$ into nonempty integer rectangles
$S=\prod_i\{a_i,\ldots,b_i\}$. Different channels may have different partitions;
within a channel the partition is shared across all spatial components $k$.
The canonical observation map has rows
\begin{equation}\label{eq:canonical}
 (Cz)_{\alpha,k,S}=V^{-1/2}|S|^{-1/2}
             \sum_{x\in S}z_{k,\alpha}(x),\qquad
 N_{\rm out}=D\sum_\alpha|\mathcal P_\alpha|.
\end{equation}
These are unrotated, fixed-weight component means. No mixed internal row or
subsequent dimension-reducing channel mixture is part of this class.
Let $\mathcal R$ be the counting orthogonal block projection replacing each
component by its arithmetic mean on its channel's rectangles. Then
\begin{equation}\label{eq:RC}
 \norm{Cz}^2=V^{-1}\norm{\mathcal Rz}^2.
\end{equation}
The \emph{actual rank} is $r(C)=\rank(C|_\UU)$, which need not be $N_{\rm out}$.
We call $C$ restoring at target $c>0$ if some finite $K\geq0$ gives
\begin{equation}\label{eq:restoring}
 V^{-1}Q_B(z)+K\norm{Cz}^2\geq cV^{-1}\norm z^2
 \quad(z\in\UU).
\end{equation}
We separately state when the same inequality holds on all of $\mathcal H_g$.
On these \emph{same} rows the counting-energy convention uses coefficient $VK$.
Arbitrary row rescaling would change the penalty question and is excluded.
The projection $\mathcal R$ is not assumed to preserve $\UU$.

\subsection{Sharp order for a fixed weak space}
Let $1\leq q\leq g$, $\Lambda\geq1$, and let $P$ be a spatially constant
rank-$q$ orthogonal projector satisfying
\begin{equation}\label{eq:coherence}
 P_{\alpha\alpha}\leq\Lambda q/g\quad(1\leq\alpha\leq g).
\end{equation}
With $c>0$ fixed put
\begin{equation}\label{eq:profile}
 a=\frac{c}{16D},\qquad H=cn^2,\qquad
 B=HI_g-(H-a)P,\qquad b=\min\{g/q,n^2\}.
\end{equation}
There is no lower bound on nonzero diagonal entries of $P$ and no prescribed
well-conditioned subset of channels.

\begin{theorem}[Fixed weak space]\label{thm:fixed}
For \eqref{eq:profile}--\eqref{eq:coherence} there are coefficient-only canonical
partitions such that
\begin{equation}\label{eq:fixedupper}
 N_{\rm out}\leq L_{D,\Lambda}(qV+gV/n^2),\qquad
 Q_B(z)/V+32cb\norm{Cz}^2\geq c\norm z^2/V
 \quad(z\in\mathcal H_g).
\end{equation}
One valid constant is
\[
 C_0=2^{24},\quad
 A_D=\frac{D2^D(17/16)}{2^{D-2}-1},\quad
 L_{D,\Lambda}=A_D+4352D+\frac{17}{4}DC_0\Lambda.
\]
Every restoring internally channel-separable linear map, even beyond canonical
rectangles, has actual rank
\begin{equation}\label{eq:ranklower}
 r\geq\max\left\{qd_0,
 \frac{gd_0}{2\Lambda(1+8Dn^2)}\right\}.
\end{equation}
This rank bound only uses the restoring inequality at zero observation; it
also holds for any finite penalty $\Phi(Cz)$ with $\Phi(0)=0$ in place of the
quadratic penalty. Consequently the minimum canonical restoring actual rank
has order $qV+gV/n^2$, with constants depending only on $D,\Lambda$.
For the particular constructed canonical map, its least quadratic coefficient
on either $\UU$ or all fields satisfies
\begin{equation}\label{eq:penaltyorder}
 \frac{3(D-1)}{16\Lambda L_{D,\Lambda}}cb
 \leq K_{\min}\leq32cb.
\end{equation}
All constants are independent of $q,g,n,c$.
\end{theorem}
Here channel-separable means a direct sum of linear maps, one acting on each
original channel's entire $D$-component field. No lower bound for internally
mixed maps is asserted. The coefficient order in \eqref{eq:penaltyorder}
concerns a fixed constructed map, not an upper bound for every map of that rank.
The construction is a finite existence prescription; neither efficient search
nor a practical condition number is established.

\subsection{Spatial variation with the same local spectra}
For even $n$ let $p(x)=\sum_i x_i\bmod2$, which is well-defined on the torus.
For $g=2,q=1,\Lambda=2$, set
\begin{equation}\label{eq:checker}
 P(x)=\begin{cases}e_1e_1^T,&p(x)=0,\\e_2e_2^T,&p(x)=1,\end{cases}
 \qquad B(x)=HI_2-(H-a)P(x),
\end{equation}
with the same $a,H$ from \eqref{eq:profile}. Only spatial constancy of $P$
has been removed; local eigenvalues, positivity, upper coherence, observed
basis, canonical weights and actual domain are unchanged.

\begin{theorem}[Orthogonal checkerboard]\label{thm:checker}
For \eqref{eq:checker} and every even $n\geq2$:
\begin{enumerate}[label=(\roman*),leftmargin=*]
\item For every $D\geq3$ a canonical map with
$N_{\rm out}=2D(n^{D-1}+1)$ restores $c$ on all fields with a finite coefficient.
\item For $D=3$ a canonical map with
$N_{\rm out}\leq2^{19}n(1+\ln n)$ restores $c$ on all fields.
\item For $D=4$ a canonical map with
$N_{\rm out}\leq2^{25}n^2$ restores $c$ on all fields.
\end{enumerate}
In (ii)--(iii), the non-singleton construction has the sufficient normalized
coefficient
\begin{equation}\label{eq:pointK}
 K=M+\frac{2M^2}{3c},\qquad M=4DH+c;
\end{equation}
the explicit singleton fallbacks use $K=c$. Actual rank on $\UU_0^2$ is at
most the corresponding full count. These are upper bounds, not optimality,
matching lower bounds, logarithmic necessity or uniform-penalty statements.
\end{theorem}
In particular a universal per-profile $\Omega(V)$ rank lower cannot be extended
from constant $P$ merely by retaining the same positive local spectra and
pointwise coherence. This does not refute a worst-case lower over a larger
class containing constant profiles, or a lower with an additional uniform
penalty restriction.

\begin{theorem}[Fixed-angle plane construction]\label{thm:angle}
Replace the two projectors in \eqref{eq:checker} by any two fixed distinct weak
lines with principal angle $\theta\in(0,\pi/2]$, in the original observed basis.
Let $\kappa=1-|\cos\theta|$ and $m_\theta=\lceil4/\sqrt\kappa\rceil$.
For every $D\geq3$ and even $n\geq2$ there is an all-field restoring canonical
map with
\[
 N_{\rm out}\leq2Dm_\theta(n^{D-1}+1)
\]
and a finite penalty. When $n\geq m_\theta$, \eqref{eq:pointK} suffices;
otherwise the singleton fallback has $K=c$.
\end{theorem}
For fixed $\theta,D$ this is $o(V)$, with constants allowed to diverge as
$\theta\downarrow0$. We prove no $n$-dependent-angle order and no extension of
the point-anchor orders in Theorem~\ref{thm:checker}(ii)--(iii) to nonorthogonal
lines. That separate question is outside the present theorem list.

\subsection{The unchanged capped budget and proof roadmap}
To fix the normalization inherited by these examples, define for positive
local gradient-stack eigenvalues $\lambda_j(x)$
\begin{equation}\label{eq:cap}
 \alpha_D=\left(\frac34\right)^D\frac D2
       \left(\frac D{D-2}\right)^{D-1},\qquad
 F_c(\mathcal A)=\sum_{x,j}\min\{1,\alpha_D(c/\lambda_j(x))^{D/2}\}.
\end{equation}
For the fixed-space family, $a$ has multiplicity $qD^2$ and $H$ has
multiplicity $(g-q)D^2$ per site. Its weak cap equals one, so
\begin{equation}\label{eq:fixedcap}
 F_c(\mathcal A)=D^2V[q+(g-q)\min\{1,\alpha_Dn^{-D}\}]
 \leq D^2(qV+\alpha_Dg).
\end{equation}
Indeed $\alpha_D(16D)^{D/2}\geq(3\sqrt D)^D>1$.
Thus $h+F_c(\mathcal A)=O_D(qV+g)$; we do not identify it with the optimal
rank expression $qV+gV/n^2$. For both checkerboards the same formula with
$q=1,g=2$ is unchanged. In particular
\[
 \alpha_3=729/128,\quad \alpha_4=81/16,\qquad
 F_c=D^2V[1+\min\{1,\alpha_Dn^{-D}\}],\quad h=2D.
\]
Every cap and multiplicity is retained.

Section~\ref{sec:finite} derives all elementary finite spectral and projection
facts. Section~\ref{sec:fixed} proves Theorem~\ref{thm:fixed}; its sole
non-elementary input is the explicitly stated primary MSS theorem.
Sections~\ref{sec:checkerproof}--\ref{sec:points} prove the actual strong-graph,
plane and point constructions. Scope, limited primary comparisons and exact
source/review context are separated in Section~\ref{sec:context}.
